\chapter*{Abstract}
\addcontentsline{toc}{chapter}{Abstract}

Global food security depends on the production stability of staple crops such as maize, rice, wheat, and soybean. However, climate variability associated with large-scale modes such as the El Niño-Southern Oscillation (ENSO) and the Madden-Julian Oscillation (MJO) can modify water availability, temperature, and the occurrence of extreme events during critical crop development stages. This thesis aims to characterize the relationship between global climate variability and cereal yields in order to identify agroclimatic regimes and climate thresholds that may support the design of agricultural parametric insurance.

The thesis is organized around two scientific papers. The first paper analyzes the spatial and temporal sensitivity of maize, rice, wheat, and soybean yields to ENSO and MJO variability, integrating climate anomalies, crop yield data, and global crop calendars. The second paper uses the outputs from the first paper to propose a methodology for identifying climate thresholds and preliminary parametric triggers for agricultural risk management and food security applications.

% TODO: complete with final methodology, datasets, study period, and original results.

\textbf{Keywords:} food security; ENSO; MJO; cereals; crop yield; parametric insurance; climate risk; crop phenology.
